\begin{helper}[Product Mass Sums to One]
Fix \(M\in\mathbb N\), \(0\le p_0\le1\), and
\(P_R,P_B\subseteq\mathrm{Fin}(M)\).  For the mass function
\[
q=\noderef{FixedSetProductModelMassFn}(p_0,P_R,P_B)
\]
on pairs \(a=(A_R,A_B)\) of subsets of \(\mathrm{Fin}(M)\), every value is
nonnegative and
\[
\sum_a q(a)=1.
\]
\end{helper}
\begin{proof}
If \(A_R\not\subseteq P_R\) or \(A_B\not\subseteq P_B\), the defining
conditional in \(\noderef{FixedSetProductModelMassFn}\) gives \(q(A_R,A_B)=0\),
so the value is nonnegative.  On the support, the value is
\[
p_0^{|A_R|}(1-p_0)^{|P_R|-|A_R|}
\cdot
p_0^{|A_B|}(1-p_0)^{|P_B|-|A_B|}.
\]
The assumptions \(0\le p_0\le1\) give \(0\le p_0\) and
\(0\le1-p_0\).  Powers of nonnegative reals are nonnegative, and products of
nonnegative reals are nonnegative.  Hence \(q(a)\ge0\) for every \(a\).

For the total mass, first split the finite sum over all pairs
\((A_R,A_B)\).  The unsupported pairs contribute \(0\), so the sum is the same
as the sum over \(A_R\subseteq P_R\) and \(A_B\subseteq P_B\).  On this support
the summand factors into a red factor depending only on \(A_R\) and a blue
factor depending only on \(A_B\).  Finite distributivity gives
\[
\sum_a q(a)=
\left(\sum_{A_R\subseteq P_R}
  p_0^{|A_R|}(1-p_0)^{|P_R|-|A_R|}\right)
\left(\sum_{A_B\subseteq P_B}
  p_0^{|A_B|}(1-p_0)^{|P_B|-|A_B|}\right).
\]
The first parenthesis is the binomial expansion of
\((p_0+(1-p_0))^{|P_R|}\), and the second is the binomial expansion of
\((p_0+(1-p_0))^{|P_B|}\).  Since \(p_0+(1-p_0)=1\), both parentheses are
\(1\), and therefore \(\sum_aq(a)=1\).
\end{proof}
